\subsection{Automorphisms of the geometric modular topos}
\label{sec:modular-topos}

We make the comparison maps explicit because the unringed topos and the
algebraic moduli object carry different specified data.  For
$k=\overline{\Q}$ and the scheme
$U=\mathcal M_{0,4,k}=\mathbf P^1_k-\{0,1,\infty\}$, there are maps
\[
 \Aut_k(U)\longrightarrow\Aut(U_{\mathrm{et}})
       \longrightarrow\Out(\pi_1^{\mathrm{et}}U).
 \tag{E.T1}
\]
The first forgets the structure sheaf and the specified action of $k$;
the second restricts to finite locally constant sheaves.  Automorphisms
of a topos are taken up to natural isomorphism.  The statement below
computes the middle group using all finite puncture groups, and the
following statement exhibits an infinite subgroup of the second
map's kernel.  This supplies actual comparison data for both meanings
of the manuscript's proposed equality.

\begin{result}{E.T1}{Proved}{An explicit group for the four-point unringed topos}
Choose a geometric generic point of $\mathbf P^1_k$.  For each finite
set $B\subset\mathbf P^1(k)$ of size at least three let
$F_B=\pi_1^{\mathrm{et}}(\mathbf P^1_k-B)$, with the filling maps
$F_{B'}\twoheadrightarrow F_B$ for $B\subset B'$.  Thus $F_B$ is free
profinite of rank $|B|-1$, with its specified peripheral procyclic
subgroups.  Put
\[
 G=\varprojlim_B F_B.
\]
For each $x\in U(k)$ let $[I_x]$ be the compatible conjugacy class of
inertia above $x$, a copy of $\widehat\Z(1)$ in $G$.  Then generic
restriction gives an isomorphism
\[
 \Aut(U_{\mathrm{et}})\ \xrightarrow{\ \sim\ }\
 \bigl\{[\alpha]\in\Out_{\mathrm{cont}}(G):
   \alpha\text{ permutes }\{[I_x]:x\in U(k)\}\bigr\}.
 \tag{E.T2}
\]
The group on the right is specified by the inverse system of free
profinite groups and filling maps just given.  Every automorphism in
it has a unique induced permutation of $U(k)$.  Moreover
$N_G(I_x)=I_x$, and generic restriction has trivial kernel.
\end{result}

\begin{proof}
Let $K=k(t)$.  Every finite separable extension of $K$ has a finite
branch set on the smooth projective curve $\mathbf P^1_k$.
Normalization therefore identifies
$G=\Gal(K^{\mathrm{sep}}/K)$ with the displayed inverse limit.
The Riemann existence comparison identifies each $F_B$ with the
completion of the punctured-sphere group
\[
 \langle c_b\ (b\in B)\mid\prod_{b\in B}c_b=1\rangle,
\]
where the order and paths defining this presentation are chosen for
each $B$ and the filling maps are the actual maps of punctured curves.
No choice of independent free bases is used to change those maps.
The comparison is \cite[Theorem 21.3]{Milne2013}; its proof and the
characteristic-zero base-extension statement were opened.
For base extension we also checked the complete curve argument of
\cite[Proposition 2.10, with Lemmas 2.3, 2.6 and 2.9]{Landesman2024}.

Fix a prolongation of the valuation at $x$.  Its decomposition group
equals its inertia group because the residue field $k$ is algebraically
closed.  Characteristic zero gives
$I_x\simeq\widehat\Z(1)$.  Its image in every $F_B$ with $x\in B$
is the full peripheral group, and the map on this group is an
isomorphism.  Such finite sets, with at least three elements, are
cofinal.  A peripheral generator is primitive in $F_B$.  Its
centralizer is its own procyclic group by
\cite[Corollary 2.3]{BuxErshovRapinchuk2011}.  An element normalizing
that group acts on its generator by a unit power; abelianization in
$\widehat\Z^{|B|-1}$ makes this power equal to one.  Its normalizer
is therefore the same procyclic group.  An element of $N_G(I_x)$
lies in its image in every such quotient.  The image in one fixed
quotient determines its unique exponent in $I_x$; subtract that
element of $I_x$ and use all the quotients to obtain the identity.
This proves $N_G(I_x)=I_x$.  For $x\ne y$ a quotient with
$x,y\in B$ and $|B|\geq3$ distinguishes their peripheral subgroups
by the primitive lines in abelianization.  Thus $[I_x]\ne[I_y]$.

Here is the finite-sheaf reconstruction that proves (E.T2).  Choose
the prolongations once and use them for all the following stalks.
A finite constructible sheaf is the following data:
\[
 \left(E,\ (A_x,\ u_x:A_x\longrightarrow E^{I_x})_{x\in U(k)}\right).
 \tag{E.T3}
\]
Here $E$ is a finite continuous $G$-set, each $A_x$ is a finite set,
and outside a finite set the inertia action on $E$ is trivial and
$u_x$ is an isomorphism.  Morphisms are an equivariant map of the
generic sets and maps of all the $A_x$ commuting with the $u_x$.
To verify this description, remove a finite set containing the
ramification of $E$ and the exceptional $u_x$.  On its complement
$E$ gives a finite locally constant sheaf.  At a removed point $x$
the stalk of its direct image is $E^{I_x}$: the punctured strict
henselization has fundamental group $I_x$.  Gluing a closed stalk
$A_x$ is exactly the map $u_x$.  More explicitly, for the open and
closed inclusions $j$ and $i$ the glued sheaf is the fibre product
\[
             j_*E\times_{i_*i^*j_*E}i_*A.
\]
This formula verifies both inverse constructions, including
morphisms.  Enlarging the finite exceptional set gives the same
sheaf, because the added $u_x$ is an isomorphism.

The coherent objects of $U_{\mathrm{et}}$ are precisely these finite
constructible sheaves.  This is the precise coherent-topos statement
in \cite[Example A.6 and Corollary A.9]{CarlsonHaineWolf2024}; we
also use its site comparison in Theorem A.8.  The quasicompact
separated etale objects give generators among them.  Consequently
an equivalence on these coherent objects preserving finite limits,
finite coproducts and covers extends uniquely to a topos
equivalence.  These operations in (E.T3) can equivalently be
checked by the preceding gluing equivalence.

A topos equivalence permutes its subterminal objects, hence the
Zariski opens.  The unique generic point of their sober space is
preserved.  Localizing along all its open neighborhoods reconstructs
$\operatorname{Spec}(K)_{\mathrm{et}}$, the topos of continuous
$G$-sets.  The restricted equivalence gives $[\alpha]\in\Out(G)$.
The closed complement of the open $U-\{x\}$ reconstructs the closed
point topos, and its gluing functor on finite generic objects is
$E\mapsto E^{I_x}$.  The latter functor is represented by the
profinite $G$-set $G/I_x$.  The pro-Yoneda correspondence shows that
its image under $\alpha$ is the functor for a point $y$ exactly when
$\alpha(I_x)$ is conjugate to $I_y$.  The distinctness proved above
makes $y$ unique.  This proves that generic restriction lands in
the right side of (E.T2).

Conversely, let $[\alpha]$ belong to the explicitly displayed group,
write its point permutation as $s$, and choose $c_x\in G$ with
$\alpha(I_x)=c_xI_{s(x)}c_x^{-1}$.  On (E.T3) define
\[
 E\longmapsto\alpha^*E,\qquad
 A_x\longmapsto A_{s(x)},\qquad
 u_x\longmapsto c_xu_{s(x)}.
 \tag{E.T4}
\]
The fixed-point target of the last map is
$E^{\alpha(I_x)}=(\alpha^*E)^{I_x}$.  The normalizer calculation
shows that a different $c_x$ acts identically on the fixed points.
Changing $\alpha$ by an inner automorphism gives a natural
isomorphism.  Finite exceptional sets remain finite; the inverse
construction uses $\alpha^{-1}$.  Thus (E.T4) extends to the required
topos equivalence.  Finally an equivalence restricting to the
identity on the generic topos must fix every point and has, at $x$,
only the possible ambiguity $N_G(I_x)/I_x=1$.  Applying (E.T3)
shows that it is naturally the identity on every coherent object
and then on the whole topos.  This proves injectivity and completes
both directions of (E.T2).
\end{proof}

\begin{result}{E.T2}{Refuted; replacement proved}{Relabelling and a visible infinite kernel}
For the unringed small etale topos of
$\mathcal M_{0,4,\overline\Q}$ the equality
$G_{0,4}=\operatorname{Im}\mathfrak S_4$ proposed on cote 143,
scan 31, fails.  The two maps in (E.T1) have the following
explicit properties:
\[
 \Aut_k(U)\simeq S_3,\qquad
 \mathfrak S_4\twoheadrightarrow S_3
 \text{ has kernel }V_4,
\]
and the second map has a subgroup isomorphic to $\Z$ in its kernel.
That subgroup also acts trivially on the completed exit category
for $\mathbf P^1$ with the three boundary points retained.
Its nonidentity elements move subterminal opens of
$U_{\mathrm{et}}$.  They are outside the image of all semilinear
algebraic automorphisms preserving $\{0,1,\infty\}$.
\end{result}

\begin{proof}
In the complex coordinate $z=x+iy$, take the diamond with center
$c=(2,2)$ and vertices
\[
 (3,2),\quad(2,3),\quad(1,2),\quad(2,1).
\]
Triangulate it by joining $c$ to these vertices.  Fix its boundary,
send $c$ to $d=(5/2,2)$, and extend affinely on the four triangles
and by the identity outside the diamond.  Call the resulting map
$h$.  Its four Jacobian determinants are $1/2,3/2,3/2,1/2$,
in cyclic order.  The four images are the corresponding triangles
from the interior point $d$, so they cover the diamond with
disjoint interiors.  This proves directly that $h$ is a
homeomorphism.  Its matrices and their inverses have rational
coefficients.  Hence it maps $\mathbf P^1(\overline\Q)$ bijectively
to itself: the real and imaginary parts of an algebraic number
are algebraic, and both affine directions preserve this property.
The exact matrices, vertex equations, positive determinants and
twenty-one iterates are checked in \path{code/e_topos_pl.py} and
its saved output.

For completeness, the algebraization of the induced site functor
uses more than the comparison for the three fixed punctures.
For every finite $B\subset\mathbf P^1(k)$, the homeomorphism gives
an equivalence of finite covers of
$\mathbf P^1(\C)-B$ and
$\mathbf P^1(\C)-h^{-1}B$.  Riemann existence and the
algebraically closed base-extension equivalence descend this to
finite etale covers over $k$.  These equivalences commute with
restriction as $B$ grows.  An arbitrary quasicompact separated
etale object over $U$ is an open subcurve of the normalization
of $\mathbf P^1$ in its finite generic algebra, component by
component.  Its finite branch set and the images of all omitted
points lie in $\mathbf P^1(k)$.  Pull back the finite branched
cover topologically, algebraize on the branch complement, and
normalize to restore the missing fibres.  Keep exactly the
specified unramified sheets and omit exactly the transported
points.  The resulting curve is etale over $U$.
The uniqueness of normalization and the fully faithful finite
cover comparison give the maps on morphisms and their
composition.  Joint surjectivity of coverings is preserved on
the geometric fibres.  Applying $h^{-1}$ gives an inverse site
equivalence.  This constructs the asserted automorphism of the
whole etale topos; equivalently it is (E.T4) for the action on
all puncture groups.

Move $c$ along $(2+t/2,2)$ for $0\leq t\leq1$ in the same four
triangles.  The point remains interior, so this is an isotopy
from the identity to $h$, fixed on neighborhoods of $0,1,\infty$.
The induced action on the punctured-sphere fundamental group is
therefore inner, and its outer class after completion is the
identity.  The isotopy is also stratified for these three points,
so it gives a natural isomorphism on the completed exit category.
On the horizontal segment from $c$ to $(3,2)$ the coordinate
map is $x\mapsto x/2+3/2$.  Thus
\[
             h^m(2+2i)=3-2^{-m}+2i\qquad(m\geq0).
\]
These points are distinct.  Pullback sends the subterminal open
$U-\{a\}$ to $U-\{h^{-1}a\}$; the powers of $h$ consequently
give distinct topos automorphisms and an injected copy of $\Z$.

Every $k$-automorphism of $U$ extends to its unique smooth
projective completion $\mathbf P^1_k$, and the extension
permutes its three missing points.  Conversely every such
M\"obius transformation restricts to $U$.  There is exactly one
for each permutation of the three points, proving $\Aut_k(U)=S_3$.
The six cross-ratio transforms are
$t,1-t,1/t,1/(1-t),(t-1)/t,t/(t-1)$; the four permutations in
the Klein group give the identity transform.  This proves the
stated relabelling map and its kernel.

Finally $h$ fixes $0,1,\infty$ and $i$ but moves $2+2i$.
A semilinear fractional-linear transformation fixing the first
three points has fractional-linear part the identity.  Its
field automorphism fixes $i$ and hence fixes $2+2i$, an immediate
contradiction.  The same argument applies to every nonzero power,
using the strictly monotone orbit and its inverse.  It proves
the asserted failure of surjectivity of the first comparison
map without discarding the algebraic structure silently.
\end{proof}

\begin{result}{E.T3}{Open}{The remaining interpretation and all-type determination}
Cote 143, scans 29--31, defines a modular topos over $\Q$ and
then asks for the automorphisms of its geometric base change.
For a small unringed etale topos the maps (E.T1) are literal
forgetful and restriction maps.  Its structure morphism to
$\operatorname{Spec}(\overline\Q)_{\mathrm{et}}\simeq\mathbf{Set}$
does not retain the structure sheaf or a specified
$\overline\Q$-algebra action.  For the algebraic moduli stack,
or its associated locally ringed topos with this algebra
action retained, those data are part of an automorphism.
At $(0,3)$ both groups of automorphism classes are trivial;
at $(0,4)$ (E.T2) and $S_3$ respectively give the precise groups
for these two constructions.  E.T4 below computes the algebraic
stack automorphism group in every stable type, E.T5 treats
the three unstable genus-zero types, and E.T6 treats unmarked
genus one.  E.T7 below computes the unringed topos group
at $(1,1)$.  The full \emph{unringed} topos groups in the
remaining stable types and a determination of which specified
structure the manuscript retains remain unsettled here.
\end{result}

\noindent\emph{Attack and exact remaining obstruction.}
The first construction was attacked by restriction to finite
locally constant sheaves.  E.T2 shows explicitly why that loses
information: its kernel already contains the displayed infinite
cyclic subgroup.  Keeping all closed-point inertia restores
injectivity at the generic point and proves E.T1.  In higher
dimension, there are positive-dimensional specializations and
their compatibility maps, so the generic group and the set of
closed-point inertia alone have not been shown to reconstruct
the topos.  The smallest higher-dimensional scheme case is
$\mathcal M_{0,5,\overline\Q}$.  The next concrete step is to
compute all specialization data for its two successive
punctured-curve fibrations and determine which automorphisms
of their generic groups preserve those data.  The
one-dimensional stack case $\mathcal M_{1,1}$ additionally
retains its generic order-two isotropy and the special
isotropy above $j=0,1728$; E.T7 computes these data and their
gluing, rather than erasing them in the scheme calculation.
The original scan images for cote 143 were not among the
supplied images.  All five supplied page transcriptions were
read, and the uncertainty about the retained structure is
preserved in the reading register.

The unstable unmarked genus-one descent problem originally
left here was subsequently computed in E.T6.  Its explicit
Jacobian comparison, cocycle, and identifications by base
$2$-automorphisms supersede that part of this open statement.
